% Inactive candidate; no active book input and no release certificate.
\subsection{Hopf surface fibers and unbounded ambient mass}
\label{subsec:hopf-surface-fibration}

The complete model space of a curve submersion can occur on a
manifold whose ambient Hermitian masses are unbounded. Blowing
up the elliptic base locus of the projection $[z_0:z_1]$ on the
scalar Hopf threefold turns this projection into a submersion
with Hopf surface fibers. The
rank-one background then retains only the curve singularities,
while the smooth potentials of the Hopf threefold still force
unbounded mass in a positive ambient background.

\begin{example}\label{fibration:unbounded-hopf-threefold}
There is a connected compact complex threefold $Z$ with a
surjective holomorphic submersion $p:Z\to\mathbb P^1$ whose
fibers are scalar Hopf surfaces. The manifold $Z$ does not have
the bounded mass property. For the pullback
$\theta=p^*\omega_{\FS}$ of the Fubini--Study form of area one,
the entire ordinary model space has the complete geodesic polar
metric of \cref{fibration:polar-model-metric}, with diameter two.
Its ideal subspace is the analytic closure, with the unique
nearest-point, $1$-Lipschitz ideal retraction. All potentials
in this background have zero extended Monge--Amp\`ere measure,
yet it contains analytic ideal models at polar distance one.
\end{example}
\begin{proof}
Put $\lambda=\mathrm e^{-1}$ and consider the vector bundle
\[
 \mathcal V=\mathcal O_{\mathbb P^1}(-1)
                  \oplus\mathcal O_{\mathbb P^1}.
\]
Let $\mathcal V^\times$ be its total space with the zero section
removed, and define
\[
 Z=\mathcal V^\times/\langle v\mapsto\lambda v\rangle.
\]
The bundle projection descends to a holomorphic map
$p:Z\to\mathbb P^1$.
A Hermitian norm on the bundle satisfies
$\|\lambda^kv\|=\lambda^k\|v\|$. The action is free and proper:
on compact subsets of $\mathcal V^\times$ the norms have positive
lower and finite upper bounds, so only finitely many powers can
move one such subset to meet another. Thus the quotient is a
complex threefold. The annulus bundle
$\{\lambda\leq\|v\|\leq1\}$ is compact over $\mathbb P^1$ and
meets every orbit, so $Z$ is compact. Bundle charts identify
$p^{-1}(U)$ with
\[
 U\times\bigl((\mathbb C^2\setminus\{0\})/
                    \langle w\mapsto\lambda w\rangle\bigr).
\]
Hence $p$ is a holomorphic submersion with connected compact
Hopf surface fibers, and $Z$ is connected.

The tautological inclusion into $\mathbb C^2$ identifies
$\mathcal V^\times$ with the incidence manifold
\[
 \left\{([s:t],z)\in\mathbb P^1\times
      (\mathbb C^3\setminus\{0\}):\ t z_0=s z_1\right\}.
\]
The last coordinate $z_2$ comes from the trivial summand. This
is the blowup of $\mathbb C^3\setminus\{0\}$ along
$L=\{z_0=z_1=0,\ z_2\ne0\}$: in the chart $s\ne0$, write $a=t/s$ and
$z_1=az_0$, while in the chart $t\ne0$ write $b=s/t$ and
$z_0=bz_1$. These are the two usual blowup charts of the ideal
$(z_0,z_1)$. The scalar contraction preserves the incidence
relation and commutes with projection to $z$. It therefore
descends to a modification
\[
 \mu:Z\longrightarrow
 H=(\mathbb C^3\setminus\{0\})/
                \langle z\mapsto\lambda z\rangle.
\]
The map is an isomorphism away from the elliptic curve
$L/\langle\lambda\rangle$. Local quotient charts show that it
is its blowup; in particular it has degree one.

Take the Gauduchon form $\omega$ and smooth potentials $V_j$ on
$H$ from \cref{thm:hopf-unbounded}. Since $Z$ is compact, a
sufficiently large multiple $\Omega$ of any Hermitian form satisfies
$\Omega-\mu^*\omega>0$. Put $W_j=V_j\circ\mu$. Then
\[
 \Omega+\ddc W_j
 =\Omega-\mu^*\omega+\mu^*(\omega+\ddc V_j)>0.
\]
Positivity of smooth mixed products and the degree-one change
of variables off the exceptional divisor give
\[
 \begin{aligned}
 \int_Z(\Omega+\ddc W_j)^3
 &\geq\int_Z\mu^*(\omega+\ddc V_j)^3\\
 &=\int_H(\omega+\ddc V_j)^3\longrightarrow+\infty.
 \end{aligned}
\]
The exceptional divisor and its image have zero measure for
smooth top-degree forms, so they do not affect the equality.
Every $W_j$ is smooth and bounded, so $v_+(\Omega)=+\infty$.
By \cref{prop:bmp-metric}, $Z$ does not have the bounded mass
property.

Apply \cref{fibration:curve-envelope-pullback,
fibration:polar-model-metric} to $p$ and the area-one form
$\omega_{\FS}$. This gives all the assertions about the full
model space. To see the stated distance explicitly, fix
$x\in\mathbb P^1$ and let $g_x$ be its normalized Green model,
with $\omega_{\FS}+\ddc g_x=\delta_x$. The models $0$ and
$g_x\circ p$ are ordinary and ideal, and the latter has a
logarithmic pole of coefficient one along the smooth reduced
fiber $p^{-1}(x)$. Thus both have analytic singularities. Their
base polar measures are $0$ and $\delta_x$, so their polar
distance is one, although their top-degree measures vanish.
\end{proof}

% This second example belongs with the Chapter1 submersion ideal foundation.
\begin{example}\label{fibration:ramified-ideal-pullback}
Exact multiplier-ideal pullback can fail at a critical fiber.
On a polydisk in $\mathbb C^n$, $n\geq1$, let
$f(z_0,\ldots,z_{n-1})=z_0^2$ and put
$v(t)=\tfrac12\log|t|^2$. At the origin,
\[
 \mathcal I(v)=\mathcal O,\qquad
 \mathcal I(v\circ f)=(z_0)
             \ne\mathcal I(v)\mathcal O.
\]
\end{example}
\begin{proof}
On the curve the coefficient is $1/2$, so the multiplier ideal
is the zeroth power of the maximal ideal. Upstairs
$v\circ f=\log|z_0|^2$ has coefficient one along the smooth
divisor $\{z_0=0\}$. Apply \cref{lem:smooth-center-multiplier-ideal}
in codimension one to get its first ideal power.
\end{proof}

This calculation concerns a fixed weight. It does not decide
whether all-weight ideal types and model envelopes still descend
across critical fibers.
